\chapter{The geometric contribution for matched spherical diagrams}
\label{app:morse-geometric-contribution}

This appendix proves a nonzero asymptotic for the geometric contribution
selected from the exact finite source, and identifies its rate with the
volume of the original filled complement. It does not assert that the
complementary contribution is smaller for every undecorated diagram.
The first part supplies the analytic estimate for the entire Fourier
period; the second supplies the original topology, peripheral words,
and metric completion. Both parts retain the finite normalization.

\section{The geometric contribution and all Fourier
charges}\label{mca:unit}

For a fixed link diagram, the exact finite formula is a coefficient in
face variables followed by a linear functional in each crossing
variable. Its local polynomial splits into three specified pieces. The
central piece is defined by the two endpoint indices of each finite
recurrence path. The other two pieces contain paths restricted to one side.
The precise degree conditions are given below. Selecting the central piece at every
crossing region gives a contribution to the original invariant; it does
not define another topological invariant.

This part proves a nonzero asymptotic for that entire geometric
contribution, including all Fourier charges, both level parities, the
derivatives at repeated spectral nodes, and the original scalar factors.
The diagram may have cups and caps between crossing regions: it need not
be a flat plat. Its white triangulation must be a simple sphere with the
specified matching and marked strand edge. The local balance
inequalities used to bound the other two pieces are not needed here.

The
\hyperref[mcg:unit]{geometric
companion} constructs the actual complement and identifies the selected
action with its hyperbolic volume after sufficiently long fixed signed
fillings. A strict estimate for the full contribution minus this
geometric term remains a separate requirement for a full-invariant
volume theorem.

\subsection{Exact assumptions and statement}\label{mca:sec-1}

Fix a finite connected Morse diagram whose actual white triangulation is
a simple sphere, with its actual matching, cyclic orders, crossing-disk
frames and marked main edge. Let the marked edge join the two protected
white faces \(u,v\). The diagram has \(R\) matching edges, \(R+2\) white
faces, and \(R\) ordinary face targets. Assume the labeled positive
convex right-angled certificate and its marked complete packing are
supplied. Set the protected curvatures to zero and retain the same
positive finite curvatures \(\rho_f\) at every occurrence of each other
face. The geometric existence of this certificate is a separate input to
the present analytic argument.

The exact general Morse source identity is the one in
\href{https://sashakolpakov.github.io/murakamis_dream/notes/research_matched_sphere_morse_source.html}{A
matched spherical triangulation has the original general Morse source},
§§5--6, equivalently the finite-source part of
\href{https://sashakolpakov.github.io/murakamis_dream/notes/actual_balanced_terminal_fan_volume.html}{Balanced
terminal decorations and volume limits beyond tangle-operation trees},
§5. In particular all physical cups, caps, the unrestricted marked row,
the integrated even-level coset and the original component/framing
factors are retained. Put \(N\ge3\), \(d=N-1\), \(q=e^{\pi i/N}\), and
\(h=\lfloor d/2\rfloor\). At a crossing region, write the homogeneous
local numerator as
\[
 H_{V_N}(x,b,c,y;\sigma)
   =\sum_{i+k+\ell+t=d}
       a_{i,k,\ell,t}(\sigma)b^i x^k y^\ell c^t .
\]
Here \(x,b,c,y\) are the four adjacent face variables, and \(\sigma\)
is the spectral parameter. The coefficients \(a_{i,k,\ell,t}\) are
the matrix coefficients prescribed by the finite source identity.
Define \(H_{V_N,\mathrm C}\) by retaining precisely the terms with
\(i\le h\) and \(t\le h\). The left piece retains \(t>h\); the right
piece retains \(i>h\). These three supports are disjoint and exhaustive,
since \(i+t\le d\). More explicitly, a recurrence path from index
\(j\) to index \(i\), with \(0\le i\le j\le d\), contributes middle
powers \(b^ic^{d-j}\). It belongs to the central piece when
\(i\le h\) and \(d-j\le h\). The word central describes these endpoint
conditions: a two-step path may skip the intervening row and still
belong to this piece. At odd levels the central zero-step path is
also included, exactly once.

Selecting the central piece at every crossing gives
\begin{equation}\label{mca:eq-1}
\begin{split}
 C_N(\boldsymbol\delta)
 &=\left[\prod_{f\ne u,v}L_f^d\right]
       \left.\prod_{e=1}^{R}
       H_{V_N,\mathrm C}(L_{x_e},L_{b_e},L_{c_e},L_{y_e};\sigma_e)
       \right|_{L_u=L_v=0},\\
 \sigma_e&=-2\cos(2\pi\delta_e).
 \end{split}
\end{equation}
The brackets extract the displayed coefficient in the ordinary face
variables. The protected variables \(L_u,L_v\) are set to zero before
extraction. Thus the definition uses the same face demands and reduced
mark as the full invariant. This is a projection of the known local
polynomial, not a new topological evaluation. Write
\(P_{\mathrm C}=(-1)^dH_{V_N,\mathrm C}\) for the positive endpoint
convention. Thus \eqref{mca:eq-1} is the raw coefficient and contains
\((-1)^{dR}\).

For fixed nonzero signed crossing powers \(n_e=\epsilon_e-2p_e\), its
exact all-central period is \begin{equation}\label{mca:eq-2}
\begin{split}
 I_{N,\mathbf n}^{\mathrm C}
 ={}&q^{d^2\sum_e n_e/4}(-i)^R
 \sum_{\boldsymbol\ell\in(\mathbb Z\setminus\{0\})^R}
       \left(\prod_e\ell_e\right)
             (-1)^{(N+1)\sum_e\ell_e}\\
 &\quad\times\int_{[-1/2,1/2]^R}
       C_N(\boldsymbol\delta)\prod_e\sin(2\pi\delta_e)
       \exp\left\{N\sum_e
         (\pi i n_e\delta_e^2+2\pi i\ell_e\delta_e)\right\}
                  d\boldsymbol\delta .
 \end{split}
\end{equation} No additional power of \(N\), modified dimension, or spectral
square-root Jacobian appears. The displayed \(q\) power is the original
unwrapped scalar \(\exp(\pi i d^2\sum n_e/(4N))\). An exterior
framing/component character multiplies the raw contraction when
converting to the specified invariant convention; it has modulus one at
complete color. The full one-cut modified convention additionally
retains its stated \(i^{N+3}/N\). These conventions are not changed by a
central projection.

There is a fixed small complex polydisk and finite holomorphic actions
\(A_N\), face Hessians \(H_N\), such that the following two source
estimates hold: \begin{equation}\label{mca:eq-H1}
C_N(\delta)=
 \frac{(-1)^{dR}N^{2R}e^{NA_N(\delta)}}
 {(2\pi N)^{(R-1)/2}\sqrt{\det H_N(\delta)}}
                     (1+O(N^{-1})),
\end{equation} uniformly and holomorphically on that polydisk, and, for every
\(\eta>0\), \begin{equation}\label{mca:eq-H2}
|C_N(\delta)|\le e^{N(\Phi-\gamma_\eta)}
 \quad\text{if }\max_e|\operatorname{Re}\delta_e|\ge\eta,
 \quad\operatorname{Re}\delta_e\in[-1/2,1/2],
 \quad|\operatorname{Im}\delta_e|\le\tau_\eta,
\end{equation} for some \(\gamma_\eta,\tau_\eta>0\) and all sufficiently large
\(N\). All fixed derivatives satisfy the corresponding bound on smaller
domains. Here \begin{equation}\label{mca:eq-3}
A_N\longrightarrow A_0,\qquad A_0(0)=\Phi,\qquad
 A_0''(0)=-2\operatorname{diag}(C_e),\quad C_e>0,
\end{equation} with every fixed derivative; \(A_N\) is separately even and real on
the real chart. The marked complete action satisfies
\(2\pi\Phi=\operatorname{Vol}(Y_D)\) under the supplied geometric
certificate and its calibrated source action.

There is a diagram-dependent \(n_0\) such that, for every fixed
independently signed vector with \(|n_e|\ge n_0\), the entire period
\eqref{mca:eq-2} has the relative asymptotic \begin{equation}\label{mca:eq-4}
\begin{split}
 I_{N,\mathbf n}^{\mathrm C}
 ={}&q^{d^2\sum_e n_e/4}(-2i)^R\sqrt{2\pi}\,N^{R+1/2}\\
 &\times
 \frac{\prod_e\sin(2\pi\delta_{N,\mathbf1,e})}
 {\sqrt{\det H_N(\boldsymbol\delta_{N,\mathbf1})}
       \sqrt{\det Q_{N,\mathbf1}}}
 e^{NS_{N,\mathbf1}}
 \left(1+O_{\mathbf n}(N^{-1})
              +O_{\mathbf n}(e^{-\kappa_{\mathbf n}N})\right).
 \end{split}
\end{equation} The quantities in \eqref{mca:eq-4} are the finite critical point,
critical value and negative Hessian of \begin{equation}\label{mca:eq-5}
\Phi_{N,\boldsymbol\ell}(\delta)
   =A_N(\delta)+\pi i\sum_e n_e\delta_e^2
                         +2\pi i\sum_e\ell_e\delta_e.
\end{equation} The face determinant root is continued from the positive complete
Hessian. The spectral root is the Gaussian branch on matrices with
positive real part. The coefficient in \eqref{mca:eq-4} tends to a
nonzero number after the explicit level factors and action are removed.
In particular the period is nonzero for all sufficiently large levels of
either parity, and \begin{equation}\label{mca:eq-6}
\lim_{N\to\infty}\frac{2\pi}{N}\log|I_{N,\mathbf n}^{\mathrm C}|
                  =2\pi\operatorname{Re}S_{0,\mathbf1}.
\end{equation} All powers are fixed before \(N\to\infty\); no bound on their ratios
is required. There is no assertion here that the full invariant differs
from \eqref{mca:eq-2} by a smaller term.

\subsection{The graph facts replacing a flat-plat
assumption}\label{mca:sec-2}

Deleting \(u,v\) from a simple triangulated sphere leaves a connected
ordinary graph. Indeed reroute a path around the first deleted vertex
along its cyclic link, then around the second along its link with the
first vertex removed. That remaining link is connected. Every surviving
edge belongs to an original shaded triangle with two positive
finite-face variables.

No matching edge can have the protected pair as its endpoints: their
unique sphere edge is the marked main edge. Consequently every crossing
circle has \(x+y>0\). At least one of its two triangles is nonmonomial.
If both outers are positive, \(xy>0\); if one outer is protected, at
most one middle can be the other protected face, so the other middle is
positive. These facts also cover the smallest simple sphere. They
require neither complete-radius balance nor a tree in the ordinary
graph.

At the actual complete packing, the circle-center angle sums make the
common face action stationary. Its Hessian is positive on the
common-scale quotient: each nonmonomial triangle detects the differences
of its finite logarithmic face radii, and the connected ordinary graph
leaves only common scaling in the kernel. The active rows are literal
\(+1,-1\) differences. Thus the quotient lattice is primitive, has
dimension \(R-1\), and has no additional angular aliases.

This is the content needed from
\href{https://sashakolpakov.github.io/murakamis_dream/notes/actual_positive_face_coefficient_saddle.html}{The
positive face coefficient has one global saddle}, §§2--4, and
\href{https://sashakolpakov.github.io/murakamis_dream/notes/actual_original_boundary_transverse_rank.html}{The
original boundary source has exactly two flat face directions}, §2.
Their proofs use the actual triangle graph and finite determinant
matrices. Morse bends introduce no new face coordinate, while the exact
source identity has already dealt with their cup/cap and mark factors.
The old global theorem's phrase ``flat plat'' is therefore removed at
this specific interface, rather than by an unproved equivalence of
presentations.

\subsection{Local analytic inputs, including the odd
normal}\label{mca:sec-3}

At each original crossing region let \(\mathcal B_{e,N}(\delta_e;L)\) be
the exactly anchored even fold interpolation. It retains the actual
maximal-channel factor, beta factors, forced products and half-index. In
particular its anchor is at \(\delta_{\min}=0\) for odd levels and
\(\delta_{\min}=1/(2N)\) for even levels. The local rigidity theorem
gives \begin{equation}\label{mca:eq-7}
P_{\mathrm C,e}(\delta_e;L)
       =N^2\mathcal B_{e,N}(\delta_e;L)
                          (1+O(e^{-cN}))
\end{equation} on a fixed complex spectral and face neighborhood, with protected
variables still exactly zero. This is
\href{https://sashakolpakov.github.io/murakamis_dream/notes/actual_cross_cofactor_interpolant_rigidity.html}{Actual
all-central cofactors and the complete-fold interpolant}, §§1--4, based on
\href{https://sashakolpakov.github.io/murakamis_dream/notes/actual_positive_face_exponential_fold_interpolation.html}{An
exponentially accurate analytic interpolation of the actual fold
factors} and the independent complex packets of the actual cofactors.

For precision, the odd central formula has the form \begin{equation}\label{mca:eq-8}
P_{\mathrm C}=G_bG_c+\omega s_ms_{m+1}xyD_bD_c,
                    \qquad\omega=-4\sin^2(\pi\delta).
\end{equation} The second term in \eqref{mca:eq-8}, including its derivative at
\(\omega=0\), is retained. The even formula is its original sum of three
central-cut products. Their central factors are analytic units near
completeness. The proof of \eqref{mca:eq-7} uses the independently
bounded complex cofactor ratios and a finite Cauchy--Lagrange formula on
both parity grids; it does not differentiate a value-only interpolation.
Thus all repeated-node normals and the odd simple node are already in
\eqref{mca:eq-7}. There is no odd regularized surrogate in
\eqref{mca:eq-1}.

The local global-angular theorem applies to each terminal cofactor with
row \(N/2+O(1)\). Two positive faces with separated circular phases give
a fixed exponential loss, uniformly over the entire real spectral
interval and a sufficiently thin complex strip. Its proof retains the
exact original node adapter, positive nodal tails, all start rows and
the polynomial cost of the finitely omitted terminal nodes. One-variable
triangles are exact monomials and impose no angular condition. This
local theorem has no global topology or Morse-order hypothesis.

Every summand of the even or odd central formula contains a cofactor
from both original triangles. Outside the common ordinary-face phase, a
bounded-length path in the connected ordinary graph has a separated
edge; one triangle at that edge supplies the exponential loss in every
central summand at its crossing region. Every other cofactor has its
polynomial endpoint majorant and the fixed strip growth
\(e^{O(N|\operatorname{Im}\delta|)}\). There are finitely many crossing
regions, so shrink the common strip to absorb the sum of those growth
costs. This proves the shared angular complement bound, on the original
common radii and with every central prefactor retained.

\subsection{The local asymptotic from the actual face
torus}\label{mca:sec-4}

Fix one ordinary face as the scale reference, and write
\(L_f=\rho_f e^{z_f}\) for the others. Define the finite phase \begin{equation}\label{mca:eq-9}
\Psi_N(z,\delta)=\frac1N\sum_e\operatorname{Log}
     \mathcal B_{e,N}(\delta_e;L(z))
       -\frac dN\sum_{f\ne u,v}\operatorname{Log}L_f(z).
\end{equation} Every finite anchor and any \(O(\log N)/N\) term remains in
\eqref{mca:eq-9}; only the explicit \(N^{2R}\) in \eqref{mca:eq-7} lies
outside. The local theorems give uniform convergence with derivatives to
\begin{equation}\label{mca:eq-10}
\Psi(z,\delta)=\sum_e
       \bigl[Q_{b_e}(\delta_e;L)+Q_{c_e}(\delta_e;L)\bigr]
                         -\sum_{f\ne u,v}\operatorname{Log}L_f.
\end{equation} The logarithms are real on the positive complete chart. The positive
face Hessian and the holomorphic implicit function theorem give the
actual nearby finite saddle \(z_N(\delta)\), uniformly on one fixed
polydisk. Put \(A_N=\Psi_N(z_N,\delta)\) and
\(H_N=\partial_z^2\Psi_N(z_N,\delta)\).

The exact homogeneous numerator degree is \(dR\), equal to the total
ordinary target. Hence the common angular variable integrates to one,
and the actual Cauchy measure on the quotient is
\((2\pi)^{-(R-1)}d\theta\). The integration lattice is primitive by the
general Morse face-potential identity; no additional multiplicity is
chosen here.

Use \eqref{mca:eq-7} near the saddle and the preceding angular gap on
its complement. Translating the original Cauchy torus by the small
vector \(z_N(\delta)\) crosses no pole: all finite faces remain in their
annuli, the numerator is polynomial and the target is its original
Laurent monomial. If the translation is localized, its side faces are
retained in the fixed angular gap. Uniform positive real Hessian and
bounded higher derivatives then give the Gaussian with coefficient one
and the determinant in \eqref{mca:eq-H1}. Shrinking the spectral
polydisk makes the absolute cost of its complex perturbation smaller
than the exponential precision in \eqref{mca:eq-7}, so the relative
error remains \(O(N^{-1})\), holomorphic with fixed derivative bounds.

Uniqueness of this finite face saddle preserves separate evenness and
reality. At zero, the mixed \(z,\delta\) derivatives vanish by evenness,
so eliminating \(z\) does not change the diagonal spectral second
derivative. Explicitly \begin{equation}\label{mca:eq-11}
C_e=\frac{\pi(T_{b_e}+T_{c_e})}{x_e+y_e}>0,
 \qquad T_t=\sqrt{(x_e+y_e)t+x_ey_e},
\end{equation} at the actual complete packing.
Section~\ref{mca:sec-2} ensures that at least one
summand is positive. Thus \eqref{mca:eq-3} holds, and on a smaller fixed
complex chart \(\operatorname{Re}A_N''\) is uniformly negative definite.
This proves \eqref{mca:eq-H1} directly for the central piece. The full-numerator
restoration in §4 of
\href{https://sashakolpakov.github.io/murakamis_dream/notes/general_full_spectral_numerator_geometric_packet.html}{The
original full spectral numerator has the marked geometric packet} is not
used.

\subsection{The gap on the entire spectral cube}\label{mca:sec-5}

The needed real-spectral estimate is a value bound, not coefficientwise
domination at arbitrary spectral parameters. The exact cleared cofactor
has the original start decomposition \begin{equation}\label{mca:eq-12}
R_j(b,\delta)=\sum_{i=0}^{j}b^i
     \prod_{r<i}\alpha_r(\delta)K_{i,j},\qquad
 \alpha_r=4\cos^2\frac{(r+1/2)\pi}{N}-4\sin^2(\pi\delta).
\end{equation} The normalized two-step tail transfer has limiting spectral radius \begin{equation}\label{mca:eq-13}
\rho(t,\delta)=(x+y)\sin(\pi t)
     +\sqrt{4xy\cos^2(\pi\delta)+(x-y)^2\sin^2(\pi t)},
 \qquad\Lambda=2\sin(\pi t)\rho.
\end{equation} The proof in
\href{https://sashakolpakov.github.io/murakamis_dream/notes/actual_global_cross_spectral_gap.html}{Global
real-spectral gap of the actual all-central source}, §§1--2, handles the
possible repeated transfer root by its Cayley--Hamilton power formula.
It uses fixed finite blocks away from \(t=0\), a uniform
\(O(\eta\log(1/\eta))\) endpoint allowance, and uniformly integrable
truncated logarithms for the exact prefix. It therefore gives a uniform
upper bound for every starting row, including in small fixed complex
face neighborhoods and a thin spectral strip. No ill-conditioned
eigenbasis or lower bound on distance to a physical node is assumed.

The exact maximization of the resulting start action is \begin{equation}\label{mca:eq-14}
\limsup\frac1N\log|R_j(b,\delta)|
       \le\max\{Q_b(\delta),\tfrac12\log b\};
\end{equation} for \(b=0\), only the first branch is present. This comes from its
unique pre-sign-change maximum and the final start endpoint; the entire
prefix integral is zero. For a nonmonomial triangle,
\(Q_b(0)>\tfrac12\log b\) when that branch is present, and
\(Q_b(\delta)<Q_b(0)\) for \(0<|\delta|\le1/2\). The derivative formula
is strictly negative whenever \(b>0\) or \(xy>0\). The sole exception is
the exact one-variable monomial, already excluded from any asserted
local gap.

Every crossing region has at least one nonmonomial triangle. Thus a
separated \(\delta_e\) supplies a strict loss relative to that crossing
region's complete rate; all other crossing regions have their complete
upper rates with arbitrarily small prescribed errors. On the small
common-face-phase region this estimate and its complex face extension
control the original integrand. On the complementary face torus use the
angular gap of Section~\ref{mca:sec-3}, which is uniform
for all real spectral parameters. Combining the two regions on the same
original radii proves \eqref{mca:eq-H2}. To make endpoint neighborhoods
precise, first use separation \(\eta/2\), then shrink the strip.
Cauchy's estimate gives the fixed derivative version. Neither the
inequality comparing left/right rates with complete radii nor a choice
of different radii at separate crossing regions occurs in this proof.

\subsection{The first packet, its phase, and charge
comparison}\label{mca:sec-6}

Write, by evenness, \begin{equation}\label{mca:eq-15}
\partial_e A_N(\delta)=-2\delta_e B_{e,N}(\delta_1^2,\ldots,\delta_R^2),
 \qquad B_{e,N}\longrightarrow B_e,\quad B_e(0)=C_e.
\end{equation} Shrink the chart until \(|B_{e,N}-C_e|\le C_e/8\). Choose first a
fixed outer real cube \(K_B=[-B,B]^R\) inside that chart. Choose the
scaled charge cutoff \(a>0\) sufficiently small relative to \(B\).
Finally choose a separate inner radius \(0<b<a/4\), to be used only for
the bounded bad charges in Section~\ref{mca:sec-7}. All
three constants are fixed from the source before the powers are chosen.
For positive real charges with \(\ell_e/|n_e|\le a\), the critical
equation is exactly \begin{equation}\label{mca:eq-16}
\delta_e=-\frac{\ell_e/n_e}
                  {1+i B_{e,N}(\delta^2)/(\pi n_e)}.
\end{equation} Uniform contraction on a small polydisk gives its unique solution for
all independent signed \(|n_e|\ge n_0\). In particular \begin{equation}\label{mca:eq-17}
|\delta_e|\le2\ell_e/|n_e|,\qquad
 \operatorname{Im}\delta_e=
 \frac{\ell_e}{\pi n_e^2}
 \frac{\operatorname{Re}B_{e,N}}
 {|1+iB_{e,N}/(\pi n_e)|^2}>0.
\end{equation} The imaginary part is positive for either sign of \(n_e\); all
coordinates and all sine factors are nonzero.

Translate the outer cube \(K_B\) to the horizontal plane through this
critical point. Reality on the real chart gives
\(\operatorname{Re}A_N(x+iv)\le A_N(0)-c|x|^2+C|v|^2\), with no linear
imaginary cost. On each actual side face, one \(|x_j|=B\), while
\(0\le v_e\le\operatorname{Im}\delta_e\). The quadratic chirps cost at
most \(CB\sum_e\ell_e/|n_e|\), and the charge terms have nonpositive
real part. The critical real value is at least
\(A_N(0)-C\sum_e(\ell_e/|n_e|)^2\). The already chosen \(a\ll B\),
followed by sufficiently large \(n_0\), therefore puts every side
strictly below that value. Thus the original local contour on \(K_B\)
has saddle coefficient \(+1\), with its side currents explicitly
smaller. The inner radius \(b\) is not used for this local saddle
contour.

On the translated horizontal plane the real Hessian is
\(\operatorname{Re}A_N''<0\). Hence the negative spectral Hessian
\(Q_{N,\ell}\) has positive real part. The Gaussian determinant branch
is fixed even for mixed signs. Applying the spectral Gaussian to
\eqref{mca:eq-H1} contributes \((2\pi/N)^{R/2}\), and the two measures
combine exactly as \begin{equation}\label{mca:eq-18}
N^{2R}(2\pi N)^{-(R-1)/2}(2\pi/N)^{R/2}
                     =\sqrt{2\pi}\,N^{R+1/2}.
\end{equation} The holomorphic relative \eqref{mca:eq-H1} error remains
\(O(N^{-1})\); the ordinary cubic Gaussian term integrates to zero,
while its square and the quartic term have this same order. Constants
may depend on the fixed charges and powers. The limiting prefactor is
nonzero by \eqref{mca:eq-17} and the determinant units.

Differentiation of the critical value, rather than comparison of
independent saddle remainders, gives \begin{equation}\label{mca:eq-19}
\partial_{\ell_e}\operatorname{Re}S_{N,\ell}
 =-\frac{2\ell_e}{n_e^2}
   \frac{\operatorname{Re}B_{e,N}}
        {|1+iB_{e,N}/(\pi n_e)|^2}.
\end{equation} Therefore, throughout the positive local charge box, \begin{equation}\label{mca:eq-20}
\operatorname{Re}S_{N,\mathbf1}-\operatorname{Re}S_{N,\ell}
       \ge\sum_e\frac{C_e}{2n_e^2}(\ell_e^2-1).
\end{equation} This comparison is valid for arbitrarily unbalanced fixed powers.
Separate evenness makes every sign choice of the same positive charge
magnitudes contribute equally: reflecting a coordinate reverses both its
sine integral and its charge factor. For the unit charges the raw sign
cancels the parity character, \((-1)^{dR}(-1)^{(N+1)R}=1\), leaving the
exact factor \((-2i)^R\) in \eqref{mca:eq-4}.

\subsection{Whole-domain currents and every other
charge}\label{mca:sec-7}

For any entire polynomial in the \(\sigma_e\), multiply it by the
original sines and quadratic chirps. In each variable the resulting
\(B\) has zero endpoint values and equal first derivatives at
\(\pm1/2\). This remains true after spectator derivatives. Three
integrations by parts give exactly \begin{equation}\label{mca:eq-21}
\int_{-1/2}^{1/2}B(t)e^{2\pi iN\ell t}\,dt
 =(2\pi iN\ell)^{-3}
 \left\{[B''(t)e^{2\pi iN\ell t}]_{-1/2}^{1/2}
       -\int_{-1/2}^{1/2}B'''(t)e^{2\pi iN\ell t}\,dt\right\}.
\end{equation} The second-derivative current is retained. With the original charge
weight this yields summable \(|\ell|^{-2}\) at fixed \(N,\mathbf n\).
Nothing about this step assumes a plat.

\eqref{mca:eq-H1} and \eqref{mca:eq-H2} first imply the shared real
supremum \(e^{N(\Phi+o(1))}\). Since the degree in each cosine variable
is at most \(d\), the elementary polynomial exterior bound gives \begin{equation}\label{mca:eq-22}
|C_N(\delta)|\le
       e^{N(\Phi+o(1))+2\pi d\sum_e|\operatorname{Im}\delta_e|}.
\end{equation} This is a bound for the whole shared coefficient. Source derivatives
are estimated by Cauchy, while spectator chirps are differentiated on
their real intervals, so no exponentially large spectator-slope cost is
inserted.

For a positive far charge in coordinate \(i\), shift its whole interval
upward by a fixed small height \(\tau\). Both endpoint sides are kept
together. Their exact difference, with spectator factors suppressed, is
\begin{equation}\label{mca:eq-23}
\begin{split}
 B_i(1/2+iv)-B_i(-1/2+iv)
 ={}&2i\sinh(2\pi v)\sinh(\pi n_iNv)\\
 &\times e^{\pi i n_iN(1/4-v^2)}
           C_N(1/2+iv,\delta_{\mathrm{other}}).
 \end{split}
\end{equation} The endpoint Fourier phase combined with the original parity is
\((-1)^{\ell_i}\), of modulus one. \eqref{mca:eq-H2} gives a fixed gap
at these separated endpoints. Applying \eqref{mca:eq-21} in all
spectator coordinates bounds the weighted vertical term by \begin{equation}\label{mca:eq-24}
C_{\mathbf n}|n_i|N\ell_i\,
    e^{N(\Phi-\gamma_e)}\prod_{j\ne i}|\ell_j|^{-2}
       \int_0^\tau v^2e^{-2\pi N(\ell_i-|n_i|/2)v}\,dv.
\end{equation} For \(\ell_i>|n_i|/2+D\), its integral is bounded by
\(C/[N^3(\ell_i-|n_i|/2)^3]\), and the resulting sum converges. The
horizontal interval uses \eqref{mca:eq-22} and has exponential rate at
most \(\Phi+o(1)+2\pi R\tau-2\pi\tau(\ell_i-|n_i|/2)\). Choose fixed
\(D\) large enough to make this uniformly smaller than \(\Phi\). The
union over any far coordinate, summed with all spectators on their
entire real intervals, has a fixed exponential gap. Negative charges
obey the same bound by reflection. No infinite series restricted to a
sharp small cube has been introduced.

The remaining charges satisfy \(|\ell_i|\le\lceil|n_i|/2+D\rceil\), a
finite set after the powers are fixed. Consider first a bounded bad
charge for which one ratio \(|\ell_i|/|n_i|>a\). Outside the inner cube
\(K_b\), \eqref{mca:eq-H2} supplies a fixed gap. Inside \(K_b\), reflect
the chosen charge to positive sign and shift only that coordinate by
\(ib/|n_i|\). Since \(0<b<a/4\), the horizontal charge loss is at least
\(2\pi b(a-b)>0\); the complex-source error is \(O(b^2/n_i^2)\). On each
actual vertical side \(|x_i|=b\), the fixed real quadratic loss
persists, and the chirp plus charge has nonpositive real part. This
proves a fixed gap for these bounded bad charges, independently of power
ratios.

For the remaining finite local charge box, with all ratios at most
\(a\), use \eqref{mca:eq-H2} outside the outer cube \(K_B\), and the
actual translated \(K_B\) contour and comparison \eqref{mca:eq-20}
inside. Every nonunit magnitude has a positive fixed-\(\mathbf n\) gap.
The unit drop from \(\Phi\) is \(O_D(\sum_e n_e^{-2})\). Enlarge the
common diagram-dependent threshold until this drop is smaller than the
previously fixed endpoint, far-charge, bounded-bad and
outer-cube-complement gaps. Also ensure \(a n_0\ge1\). All complementary
contributions are then exponentially smaller than the nonzero unit
packet. This proves \eqref{mca:eq-4}--\eqref{mca:eq-6} for the whole
all-central period.

\subsection{The geometric interface and the precise remaining
problem}\label{mca:sec-8}

The analytic proof retains the untwisted selected action \(A_0\) and the
original half twists exactly once. Put \begin{equation}\label{mca:eq-25}
A_{\epsilon}=A_0+\pi i\sum_e\epsilon_e\delta_e^2,
 \qquad U_{\epsilon,e}=\partial_{\delta_e}A_{\epsilon},
 \qquad V_e=-4\pi i\delta_e.
\end{equation} The actual stationary equation of \eqref{mca:eq-5}, with
\(n_e=\epsilon_e-2p_e\), is \begin{equation}\label{mca:eq-26}
U_{\epsilon,e}+p_eV_e=-2\pi i\ell_e.
\end{equation} Once the original ordered Morse diagram has the marked
action/primitive collar adapter and smooth long-filling completion, its
regulator boundary term cancels the Gaussian terms at \(\ell_e=1\). Then
\(2\pi\operatorname{Re}S_{0,\mathbf1}\) is the actual filled exterior's
hyperbolic volume. This identifies the asymptotic of the sector
\eqref{mca:eq-2}, not automatically that of the full invariant.

The independent issue remains a bound or cancellation statement for the
entire original left/right remainder. Complete-radius balance is one
sufficient mechanism in the terminal-fan theorem. The present proof does
not assume it and supplies no replacement. In particular a nonzero
sector with the correct geometric rate can still be canceled by another
sector unless that remainder is controlled.

\subsection{Proof dependencies and scope}\label{mca:sec-9}

The following table records the inputs and the hypotheses used in each
part of the proof.

\begingroup
\begin{longtable}[]{@{}
  >{\raggedright\arraybackslash}p{(\linewidth - 4\tabcolsep) * \real{0.3333}}
  >{\raggedright\arraybackslash}p{(\linewidth - 4\tabcolsep) * \real{0.3333}}
  >{\raggedright\arraybackslash}p{(\linewidth - 4\tabcolsep) * \real{0.3333}}@{}}
\toprule\noalign{}
\begin{minipage}[b]{\linewidth}\raggedright
Input
\end{minipage} & \begin{minipage}[b]{\linewidth}\raggedright
Part used here
\end{minipage} & \begin{minipage}[b]{\linewidth}\raggedright
Dependence on a flat plat or balance
\end{minipage} \\
\midrule\noalign{}
\endhead
\bottomrule\noalign{}
\endlastfoot
\href{https://sashakolpakov.github.io/murakamis_dream/notes/research_matched_sphere_morse_source.html}{A
matched spherical triangulation has the original general Morse source}
§§5--6 & Exact coefficient, primitive integer face grid, both parity
cosets, mark and whole Hermite functional & Proved for the actual
general Morse construction; no balance \\
\href{https://sashakolpakov.github.io/murakamis_dream/notes/actual_positive_face_exponential_fold_interpolation.html}{An
exponentially accurate analytic interpolation of the actual fold
factors} §§1--6 & Exact finite anchor and exponential interpolation in
fixed complex face neighborhoods & Local actual triangles and
protected-zero patterns only \\
\href{https://sashakolpakov.github.io/murakamis_dream/notes/actual_cross_cofactor_interpolant_rigidity.html}{Actual
all-central cofactors and the complete-fold interpolant} §§1--4 & Actual
off-node central polynomial, including odd derivative term & Independent
actual cofactor growth plus interpolation; local only \\
\href{https://sashakolpakov.github.io/murakamis_dream/notes/actual_cross_cofactor_global_angular_gap.html}{A
global angular gap for the actual cleared cross cofactors} §§1--4 &
Whole real-spectral angular gap and omitted-terminal-node interpolation
& Local positive finite matrices and physical recurrences only \\
\href{https://sashakolpakov.github.io/murakamis_dream/notes/actual_positive_face_coefficient_saddle.html}{The
positive face coefficient has one global saddle} §§2--4 & Positive
Hessian, primitive phase constraints and coefficient Gaussian & Actual
connected ordinary triangle graph; no flat gap forest \\
\href{https://sashakolpakov.github.io/murakamis_dream/notes/actual_original_boundary_transverse_rank.html}{The
original boundary source has exactly two flat face directions} §2 &
Ordinary graph connectivity after the protected pair is removed & Simple
sphere and marked main edge \\
\href{https://sashakolpakov.github.io/murakamis_dream/notes/actual_global_cross_spectral_gap.html}{Global
real-spectral gap of the actual all-central source} §§1--4 & Uniform
start-row/tail estimates, strict real-spectral decrease and shared strip
gap & Local recurrences plus the same ordinary graph and complete
radii \\
\href{https://sashakolpakov.github.io/murakamis_dream/notes/general_full_spectral_numerator_geometric_packet.html}{The
original full spectral numerator has the marked geometric packet} §§1--3
& Finite phase, exact quotient measure and actual Cauchy contour &
Reproved above with the Morse source; its full restoration §4 is
unused \\
\href{https://sashakolpakov.github.io/murakamis_dream/notes/general_localized_unit_charge_filling_period.html}{The
original localized spectral period has a nonzero unit-charge filling
contribution} §§2--5 & Fixed independent signed saddle, actual side
currents, nonzero first packet & H1 and exact scalar/evenness only; the
old flat-plat opening is not used \\
\href{https://sashakolpakov.github.io/murakamis_dream/notes/general_all_cross_long_filling_volume.html}{The
entire original all-central Fourier period from its local packet and
global spectral gap} §§2--5 & Whole-domain charge tails and finite
bad-charge separation & H1, H2 and entire polynomial degree only \\
\href{https://sashakolpakov.github.io/murakamis_dream/notes/actual_balanced_terminal_fan_volume.html}{Balanced
terminal decorations and volume limits beyond tangle-operation trees}
§§6--8 & Explicit existing Morse version of these same analytic steps &
Balance is used in §9 for the omitted full remainder, not in §§6--8 \\
\end{longtable}
\endgroup

No new analytic hypothesis is imposed on a certified simple-sphere Morse
source beyond these literal inputs. The original presentation's physical
normalization and its marked geometric applicability must still be
supplied, as stated in Section~\ref{mca:sec-1}. There is
no full-invariant volume theorem for general unchanged matched spheres
in this appendix.


\section{The original complement and its geometric
contribution}\label{mcg:unit}

A simple triangulated sphere with a perfect matching of its cubic dual
determines an actual augmented link diagram when the cyclic orders,
strand pairings and crossing disks are retained. This section follows that
diagram through its exact finite source, its geometric critical point,
and the filling of its original auxiliary cusps.

For independently sufficiently long fixed signed twists, the geometric
contribution to the original invariant is nonzero and has the
exponential rate of the actual filled hyperbolic volume. The proof keeps
the primitive peripheral words, the finite scalar factors and both
parities of the quantum level. No terminal tangles are inserted, and no
local balance assumption is imposed.

The central piece is the part of each local numerator defined by the
degree conditions in Section~\ref{mca:sec-1}. Equivalently, it consists
of the central recurrence paths of each
original homogeneous numerator. Selecting it at every crossing region is
an algebraic decomposition of the original evaluation, not a new
invariant. The
\hyperref[mca:unit]{analytic
companion} proves its full Fourier asymptotic. To conclude the same
limit for the full invariant of an undecorated diagram, one must
additionally control the remaining pieces.

\subsection{Precise diagram and filling scope}\label{mcg:sec-1}

Let \(\mathcal T\) be a finite simple triangulated oriented sphere with
at least four vertices. Choose a perfect matching of its cubic dual. The
two triangles in each matched pair form a quadrilateral with four
distinct white vertices. Record a cyclic order \((x_e,b_e,y_e,c_e)\);
its common diagonal is \(x_ey_e\), its input middle is \(b_e\), and its
output middle is \(c_e\). There are \(R\) matched pairs and \(R+2\)
white vertices. Mark a nonmatching edge \(uv\). It is a strand edge in
the link diagram.

There are two equivalent permitted inputs:

\begin{itemize}
\item
  an actual connected augmented diagram in \(S^3\), whose projection and
  crossing-disk cut has exactly these labels and orders; or
\item
  the literal inverse-cut construction in
  section~\ref{mcg:sec-2} below.
\end{itemize}

Fix the retained half twist \(\epsilon_e\in\{-1,0,1\}\) and its sign at
every disk. Its auxiliary peripheral basis \((\mu_e,\lambda_e)\) is the
original disk-framed primitive basis. The longitude \(\lambda_e\) is
represented by the ordered two-head product \(P_e=B_eA_e\). Integer
filling coordinate \(p_e\) means filling the slope
\(\mu_e+p_e\lambda_e\); the main crossing power is \begin{equation}\label{mcg:eq-1}
n_e=\epsilon_e-2p_e.
\end{equation} The original outside terminal matching and the even/odd pairing at
each disk determine the main components. When their terminal graph is
connected, denote the filled knot by \(K_{\boldsymbol p}\). We keep this
original parity test. We do not change the parity vector to obtain a
different knot.

A supplied labeled positive right-angled cut certificate is sufficient
for everything below. In this simple-sphere class such a certificate
also follows from the constructive packing theorem in
\href{https://sashakolpakov.github.io/murakamis_dream/notes/actual_balanced_terminal_fan_volume.html}{the
balanced-fan note, section 2}. That section constructs the packing
\textbf{before any fan is inserted}. Only that initial geometric
construction is used here.

\textbf{Theorem.} For each fixed diagram and the recorded choices above
there is a number \(T_D\) such that, for every fixed integer vector
\(\boldsymbol p\) with \(\min_e|n_e|>T_D\), independently in sign and
relative size, the following assertions hold.

\begin{enumerate}
\def\labelenumi{\arabic{enumi}.}
\item
  Filling the specified auxiliary components gives a complete
  finite-volume hyperbolic metric on the original filled link exterior,
  with every main cusp unfilled. If the recorded component test is a
  knot, it is \(S^3\setminus K_{\boldsymbol p}\).
\item
  The all-central term of the original complete-color Morse contraction
  has a nonzero relative first-charge asymptotic, at both parities of
  the level \(N\). Its selected critical action \(S_{\boldsymbol 1}\)
  satisfies \begin{equation}\label{mcg:eq-2}
2\pi\operatorname{Re}S_{\boldsymbol 1}
         =\operatorname{Vol}(S^3\setminus K_{\boldsymbol p})
\end{equation} in the knot case, and the analogous original link-complement
  identity otherwise.
\item
  In particular, for that all-central contraction
  \(\mathcal I_{N,\mathrm C}\), \begin{equation}\label{mcg:eq-3}
\lim_{N\to\infty}\frac{2\pi}{N}
           \log|\mathcal I_{N,\mathrm C}|
         =\operatorname{Vol}(S^3\setminus K_{\boldsymbol p}).
\end{equation} The exact raw phase is retained in
  section~\ref{mcg:sec-3}. An ordinary zero-framing
  phase does not change \eqref{mcg:eq-3}.
\end{enumerate}

The diagram, its matching, its signs, and all powers are fixed before
\(N\). The threshold can depend on this entire finite diagram. There is
no short-filling assertion and no uniformity in a growing sphere. The
theorem does not assert \eqref{mcg:eq-3} for the full invariant:
that requires a bound on the remaining sectors.

\subsection{Why the manifold and its complete anchor are the original
ones}\label{mcg:sec-2}

Take disjoint small neighborhoods of the matched edges of the cubic
dual. Each is a four-terminal disk. In cyclic terminal order
\(0,1,2,3\), the two shaded halves group \((1,2)\) and \((3,0)\). Draw
the physical untwisted strands instead with pairs \((0,1),(2,3)\), and
connect outside terminals along the nonmatching dual edges. In the
corresponding disjoint three-balls put the perpendicular crossing
circles and their twice-punctured disks. This constructs a link in
\(S^3\), not a manifold assigned to an abstract tensor graph. Retained
odd half twists use the physical pairing \((0,2),(1,3)\), with their
original signed collar paths.

The two main pairings and the shaded grouping must not be confused. The
graph of main components is the union of the outside terminal matching
and these physical pairings. Every terminal has degree two. Its
connectedness is exactly the knot test, and adding full twists does not
change it.

Cut along the projection sphere. In each projection ball a crossing
circle semicircle is boundary parallel by its original half disk.
Drilling that semicircle and opening that same half disk cancels the
introduced handle. These operations occur in disjoint collars, so
opening all of them leaves one ball in each projection copy. The
boundary cellulation is the prescribed medial sphere. Thus cycles in the
matched graph create no unidentified handle or extra gluing.

Let \(Q\) be its positive ideal right-angled cut polyhedron. The even
shaded pairing fixes the crossing-circle vertex and sends
\(C_1\mapsto C_0,\ C_2\mapsto C_3\). An odd pairing exchanges the
projection copies and the two punctures, using the literal signed
half-twist collar. The two odd signs have the same corner permutation
but different primitive paths. White faces pair by their original
projection labels.

At every ideal vertex, the link of \(Q\) is a rectangle. Normalize its
opposite shaded sides to length one. Each shaded triangle map then
preserves the canonical horocycles at all three corners, including the
odd corner permutation. Each alternating white/shaded edge cycle has
four right-angle pieces. It fixes both ideal endpoints, preserves the
chosen horospheres there, and rotates by \(2\pi\), so its whole holonomy
is the identity. The ideal links are the actual link's torus links and
have no cone or dilation defect. Their horoball ends are complete. Hence
\begin{equation}\label{mcg:eq-4}
\operatorname{Vol}(Y_D)=2\operatorname{Vol}(Q),
       \qquad \Phi_D=\frac{\operatorname{Vol}(Y_D)}{2\pi}.
\end{equation} Here \(Y_D\) is the same augmented exterior, with the recorded half
twists and peripheral bases.

These are the local cut and gluing arguments in
\href{https://sashakolpakov.github.io/murakamis_dream/notes/general_augmented_plat_matching_geometry.html}{the
matching-geometry proof, sections 4--6}. They use a sphere projection,
disjoint half disks, and the prescribed corner maps. Moving cups and
caps to a flat plat is not a premise. The inverse construction and its
general Morse realization are spelled out in
\href{https://sashakolpakov.github.io/murakamis_dream/notes/research_matched_sphere_morse_source.html}{the
matched-sphere source, sections 1--4}.

For completeness, the initial simple-sphere packing theorem does not
require the absence of separating triangles. Remove one triangular face
and pool its three boundary angle demands. Strict Hall inequalities for
every nonempty proper destination subset give a positive allocation of
angles. The strictly concave calibrated triangle entropy has an interior
maximizer, whose multipliers give positive radii. The center triangles
have zero interior rotational and translational holonomy. Their boundary
develops once around a convex triangle, so the degree argument makes the
development injective. Vertex-disk sectors stop before opposite center
edges; this proves disjointness, including across separating triangles.
The contact circles and outward halfspaces therefore give the labeled
strict support certificate. Sending the marked contact to infinity gives
the two protected parallel lines. No radius-balance inequality enters
this construction.

\subsection{The exact source and the phases which remain outside the
action}\label{mcg:sec-3}

Put \(d=N-1\), \(q=e^{\pi i/N}\). In the recorded four-leg order the
full local homogeneous numerator is
\(H_e(L_{x_e},L_{b_e},L_{c_e},L_{y_e})\). The original reduced Morse
contraction is \begin{equation}\label{mcg:eq-5}
[\prod_{f\ne u,v}L_f^d]\,
       \prod_e H_e(L_{x_e},L_{b_e},L_{c_e},L_{y_e})
             \big|_{L_u=L_v=0}.
\end{equation} This statement includes the complete root Hermite normals and the
unrestricted physical mark.

Here is the part of its proof that matters for changing the Morse height
function. Each full coupon contributes \(N^{-1}H_e/(b_ec_e)^{d/2}\). If
\(M_f\) counts middle passages and \(T_f\) counts outside minus inside
turns, the actual boundary turning number gives \begin{equation}\label{mcg:eq-6}
M_f-T_f=2 \quad(f\ne u),\qquad
       M_u-T_u=-2.
\end{equation} The ordinary caps and cups give the same powers as these winding
numbers; replacing the exposed cap by the original reduced mark
contributes \(N^{-1}(L_v/L_u)^d\). Thus the full summand is
\(N^{-R-1}\prod H_e/\prod_{f\ne u,v}L_f^d\). The primitive
face-potential lattice has \(R+1\) independent \(N\)-point grids after
the exterior gauge is fixed. Their count cancels \(N^{-R-1}\). Exact
endpoint support is imposed before the root reduction, so no coefficient
is aliased into a wrong endpoint. This proves \eqref{mcg:eq-5} with its
original phases at arbitrary Morse height and both level parities.

Replace each homogeneous factor in \eqref{mcg:eq-5} by its central
cofactor factor and retain the physical spectral trace. Write
\(C_N^{\mathrm{raw}}(\boldsymbol\delta)\) for the resulting entire
spectral numerator, with each node \(-2\cos(2\pi\delta_e)\). It is even
separately in every coordinate, has degree at most \(d\) in each node,
and satisfies \(C_N^{\mathrm{raw}}=(-1)^{dR}C_N^+\) in the positive
cofactor normalization. The exact original filling period is \begin{equation}\label{mcg:eq-7}
\begin{split}
\mathcal I_{N,\mathrm C}
={}&q^{d^2\sum_e n_e/4}(-i)^R
 \sum_{\boldsymbol\ell\in(\mathbb Z\setminus\{0\})^R}
 \Bigl(\prod_e\ell_e\Bigr)(-1)^{(N+1)\sum_e\ell_e}\\
&\quad\cdot
 \int_{[-1/2,1/2]^R}
  \Bigl(\prod_e\sin(2\pi\delta_e)\Bigr)
  C_N^{\mathrm{raw}}(\boldsymbol\delta)
  \exp\!\left[
   N\sum_e\{\pi i n_e\delta_e^2+2\pi i\ell_e\delta_e\}
       \right]\,d\boldsymbol\delta .
\end{split}
\end{equation} The sum is the original trace identity with its endpoint currents; it
is not defined by an unqualified termwise Poisson rearrangement. The
odd-level simple node and the double-node normals are included. The same
finite proof is
\href{https://sashakolpakov.github.io/murakamis_dream/notes/research_matched_sphere_morse_source.html}{the
matched-sphere source, sections 5--7}.

No operation in \eqref{mcg:eq-5}--\eqref{mcg:eq-7} changes a disk, a
strand pairing, or a peripheral curve. In particular the all-central
source for the original sphere exists before a fan construction, and its
ordinary-face targets are still \(d\).

\subsection{One marked circle pattern on the whole
sphere}\label{mcg:sec-4}

Normalize the complete protected white disks to
\(\operatorname{Im}z\le0\) and \(\operatorname{Im}z\ge1\). For each
finite face use center \((X_f,Y_f)\) and radius \(r_f\). At a strand
edge set \(I_{fg}=1\); at matching edge \(e\) set
\(I_{fg}=\cos(2\pi\delta_e)\). The equations are \begin{equation}\label{mcg:eq-8}
\begin{aligned}
(X_f-X_g)^2+(Y_f-Y_g)^2
 &=r_f^2+r_g^2+2I_{fg}r_fr_g,\\
Y_f&=r_fI_{fu},\qquad
1-Y_f=r_fI_{fv},
\end{aligned}
\end{equation} where the last two equations are imposed at the appropriate marked
line incidences. Fix one horizontal translation. The marked line-line
edge is automatic; the remaining system has \(3R-1\) equations and
variables.

Its complete Jacobian is invertible. The infinitesimal angle sums and
the positive finite-face source Hessian first force all logarithmic
radius variations to be one common scale. Each of the two triangles
adjacent to \(uv\) contains a finite circle tangent to both fixed lines.
Its radius is \(1/2\), so its scale variation is zero. All radius
variations are therefore zero.

With radii fixed, each finite center triangle has one infinitesimal
Euclidean motion. A one-line triangle has only a horizontal translation.
Removing the two triangles at \(uv\) leaves a triangulated disk with
connected triangle adjacency, containing all finite vertices and at
least one one-line triangle. Two adjacent finite triangles agree at two
distinct centers; two adjacent one-line triangles agree at their finite
center. These comparisons propagate one horizontal translation through
the disk, and the normalization removes it. The kernel is zero.

The holomorphic implicit-function theorem now gives one circle pattern
on a fixed polydisk in all \(\delta_e\). The variables \(X,Y,r\) are
complexified independently, without conjugating parameters. The solution
is real on real parameters and even in each \(\delta_e\).

For one source triangle put \(x=1/r_x,\ b=1/r_b,\ y=1/r_y\), \begin{equation}\label{mcg:eq-9}
\begin{aligned}
A&=x+y,& s&=\sin(\pi\delta),\\
D&=A^2-4xy s^2,&
T_b&=\sqrt{bA+xy+b^2s^2},\\
E_b&=A+2bs^2+2sT_b .
\end{aligned}
\end{equation} Use positive complete roots and their analytic continuations. The
original determinant, beta, and forced-product rate \(Q_b\) obeys \begin{equation}\label{mcg:eq-10}
\partial_\delta Q_b=-\operatorname{Log}(E_b/\sqrt D),
\qquad
\partial_{\log L_f}Q_b=\frac{\theta_f}{2\pi},
\qquad Q_b(0)=F(x,b,y).
\end{equation} The angles are those of \eqref{mcg:eq-8}, with their marked-line
limits. Thus the developed angles sum to \(2\pi\) exactly when the
finite-face source is stationary. Uniqueness of the finite-face critical
point identifies these curvatures with the limiting source saddle.

The common strip also handles a triangle containing both protected
lines. If its finite radius is \(r\), its strip width is
\(r(1+\cos(2\pi\delta))\). This shared width cannot be chosen
independently at two triangles. The corresponding repeated scalar root
does not remove the selected eigenflag: the noncentral complete
parabolic has a nonzero off-diagonal matrix entry, and its eigenline for
\begin{equation}\label{mcg:eq-11}
J=-e^{-2\pi i\delta}
\end{equation} continues holomorphically, including its collision at zero.

This is the complete rank and face-identification mechanism in
\href{https://sashakolpakov.github.io/murakamis_dream/notes/actual_marked_deformed_circle_pattern.html}{the
marked-pattern proof}. No part of it refers to the number of strands at
a Morse slice. Simplicity is important: in particular a matching edge
cannot join the two protected faces, since their unique common edge is
the marked strand edge. Hence \(A>0\) in every local formula.

\subsection{Full group relations and the original primitive peripheral
words}\label{mcg:sec-5}

Choose the two holomorphic white-reflection matrices
\(\mathcal R_f^\pm\), with \(\mathcal R_f^-\mathcal R_f^+=I\), from the
one pattern. At a regular Morse slice with regions \(F_0,\ldots,F_k\),
the original heads and prefixes are \begin{equation}\label{mcg:eq-12}
H_j=\mathcal R_{F_j}^+\mathcal R_{F_{j-1}}^-,
 \qquad
 H_j\cdots H_1=\mathcal R_{F_j}^+\mathcal R_{F_0}^- .
\end{equation} Their marked lifts give the original parabolic main meridians. Shared
physical wire ends have the same adjacent region labels. At a cup or cap
the consecutive regions are \(F,G,F\), so the inserted ordered pair has
product \(I\), and all later prefixes are unchanged. This proves the
relations also when \(k\) changes.

At one crossing disk the ordered input pair is \begin{equation}\label{mcg:eq-13}
A_{\mathrm{in}}=\mathcal R_b\mathcal R_x,\qquad
 B_{\mathrm{in}}=\mathcal R_y\mathcal R_b,\qquad
 P=B_{\mathrm{in}}A_{\mathrm{in}}=\mathcal R_y\mathcal R_x,
\end{equation} and the output pair replaces \(b\) by \(c\). The labeled contact map
of the disk, taking \(C_1\mapsto C_0,\ C_2\mapsto C_3\), is \begin{equation}\label{mcg:eq-14}
R_0=\frac{E_bE_c}{D},\qquad
 C=I+\frac{R_0-1}{J-J^{-1}}(P-J^{-1}I).
\end{equation} It sends both input generators to their respective output generators
by simultaneous conjugation and commutes with \(P\).

The primitive normalization in \eqref{mcg:eq-14} is geometric. Sending
the two outer intersection points to \(0,\infty\) makes the two shaded
contact circles concentric, with radii \(\sqrt D/E_b\) and
\(E_c/\sqrt D\). Their ordered corner map is multiplication by \(R_0\);
the product of their reflections would instead give \(R_0^2\). The
original threading word \([2,1,1,2]\) follows the inverse of the first
map. At zero, \eqref{mcg:eq-14} retains the nontrivial parabolic
divided-difference limit. Setting \(C=I\) merely because its eigenratio
is one would give the wrong group representation.

Marked lines are treated by one common Möbius regularization of the
whole pattern. The input/output side signs are calibrated by the same
complete cusp rectangle. Nonvanishing middle contributions have the
original opposite signs; an exceptional two-line triangle has \(T_b=0\)
and contributes ratio one. Its actual eigenflag still comes from
\eqref{mcg:eq-8}. This proves \eqref{mcg:eq-14} with the correct labels
at every zero pattern, not just at positive local curvatures.

The positive crossing acts by \((A,B)\mapsto(A^{-1}BA,A)\), represented
by \(P^{-1}-I\). For \(\epsilon\) equal to \(0,1,-1\), put \begin{equation}\label{mcg:eq-15}
\begin{aligned}
 U&=-\operatorname{Log}R_0-\epsilon\operatorname{Log}(-J),\\
 V&=2\operatorname{Log}(-J)=-4\pi i\delta,\qquad
 m=e^{U/2}.
\end{aligned}
\end{equation} All logarithms are the germs vanishing at the complete point. The
actual determinant-one auxiliary meridian is \begin{equation}\label{mcg:eq-16}
{\mathsf M}
 =mR_0(J^{-1}-1)^{-\epsilon}
        C^{-1}(P^{-1}-I)^\epsilon.
\end{equation} Its selected eigenvalue is \(m\), and
\(\det{\mathsf M}=m^2R_0(-J)^\epsilon=1\). The factor \(J^{-1}-1\) is a
unit near completeness, including for negative \(\epsilon\). Formula
\eqref{mcg:eq-16}, not just \eqref{mcg:eq-15}, verifies the original
threading and signed Hurwitz relations.

Equations \eqref{mcg:eq-12}--\eqref{mcg:eq-16} cover every shared wire,
cup, cap, drilled box, and half twist in the actual Morse van Kampen
presentation. There is no extra consistency equation at a cycle of the
matched graph: all heads came from the same global reflections and all
ordered box maps use those heads. The resulting family
\(\rho_{\boldsymbol\delta}\) is an actual framed representation of
\(\pi_1(Y_D)\). At zero it is the complete cut representation. Every
main meridian remains noncentral parabolic; its commuting longitude has
multiplier one. Thus every unfilled main peripheral group is parabolic.

These are full matrix and based-word assertions from
\href{https://sashakolpakov.github.io/murakamis_dream/notes/actual_deformed_circle_primitive_collar.html}{the
primitive collar proof} and
\href{https://sashakolpakov.github.io/murakamis_dream/notes/actual_complete_core_geometric_deformation.html}{the
global assembly, section 2}. The latter's displayed \(2k\) head count is
replaced event by event by the actual \(k(t)\) in \eqref{mcg:eq-12}. No
other part of its proof changes.

\subsection{Absolute action and the first-charge filling
equations}\label{mcg:sec-6}

Let \(A_0\) be the pure limiting face-critical value \begin{equation}\label{mcg:eq-17}
A_0(\boldsymbol\delta)=
 \operatorname{crit}_{L}
 \left\{\sum_e(Q_{b_e}+Q_{c_e})
                -\sum_{f\ne u,v}\operatorname{Log}L_f\right\}.
\end{equation} It is separately even, real on real parameters, and anchored by
\(A_0(0)=\Phi_D\). Put \[
 A_\epsilon=A_0+\pi i\sum_e\epsilon_e\delta_e^2.
\] Envelope differentiation of \eqref{mcg:eq-10}, retaining
\eqref{mcg:eq-15}, gives \begin{equation}\label{mcg:eq-18}
dA_\epsilon=\sum_e U_e\,d\delta_e,\qquad
       d(2\pi i A_\epsilon)=-\frac12\sum_e U_e\,dV_e.
\end{equation} This derivative identifies the actual primitive variables, including
the half-twist shear; it is stronger than a cusp-Hessian match.

The marked regulator boundary formula, in squared multiplier logs, is \begin{equation}\label{mcg:eq-19}
d\operatorname{Vol}(\rho)
   =\frac12\sum_e
       \{\operatorname{Re}V_e\,d\operatorname{Im}U_e
           -\operatorname{Re}U_e\,d\operatorname{Im}V_e\}.
\end{equation} The main terms vanish. Equations \eqref{mcg:eq-18},\eqref{mcg:eq-19}
have the same differential after the boundary correction, and
\eqref{mcg:eq-4} fixes the real constant. Therefore on this one selected
branch \begin{equation}\label{mcg:eq-20}
2\pi\operatorname{Re}A_\epsilon
   =\operatorname{Vol}(\rho)
        -\frac12\sum_e
                \operatorname{Im}U_e\,\operatorname{Re}V_e.
\end{equation} The logarithms are not reset during continuation.

For the positive unit charge in \eqref{mcg:eq-7}, the limiting phase is
\begin{equation}\label{mcg:eq-21}
S_{\boldsymbol1}(\boldsymbol\delta)
    =A_0(\boldsymbol\delta)
       +\pi i\sum_e n_e\delta_e^2+2\pi i\sum_e\delta_e
    =A_\epsilon-2\pi i\sum_e p_e\delta_e^2
                          +2\pi i\sum_e\delta_e .
\end{equation} Its stationarity equation is exactly \begin{equation}\label{mcg:eq-22}
U_e+p_eV_e=-2\pi i.
\end{equation} For each fixed diagram, existence and uniqueness near zero hold for
all independently long signed powers. Indeed \begin{equation}\label{mcg:eq-23}
A_0''(0)=-2\operatorname{diag}(C_e),\qquad
 C_e=\frac{\pi(T_{b_e}(0)+T_{c_e}(0))}{x_e+y_e}>0,
\end{equation} and separate evenness makes \(\partial_e A_0=\delta_e B_e\). After
multiplying the critical equations by \(1/n_e\), the derivative in
\(\boldsymbol\delta\) at reciprocal slopes zero is the invertible
diagonal \(2\pi i I\). The several-parameter implicit-function theorem
gives \(\delta_e=O_D(1/|n_e|)\). It allows independently unbalanced
signed \(n_e\); it does not compare a minimum error rate with a sum of
unrelated filling losses.

At a critical point, substitute \eqref{mcg:eq-22} into
\eqref{mcg:eq-20}. Since \(V=-4\pi i\delta\), the real part of the two
explicit added terms in \eqref{mcg:eq-21} is \begin{equation}\label{mcg:eq-24}
\operatorname{Re}(-2\pi i p\delta^2+2\pi i\delta)
       =\frac{1}{4\pi}\operatorname{Im}U\,\operatorname{Re}V .
\end{equation} The boundary corrections cancel exactly: \begin{equation}\label{mcg:eq-25}
2\pi\operatorname{Re}S_{\boldsymbol1}
                        =\operatorname{Vol}(\rho).
\end{equation} There is no omitted complex endpoint term. To replace the right side
by the hyperbolic volume of the filled topology one still needs the
following metric argument.

\subsection{Completing the specified hyperbolic
filling}\label{mcg:sec-7}

The construction in
\href{https://sashakolpakov.github.io/murakamis_dream/notes/general_source_regulator_hyperbolic_equality.html}{the
general metric completion theorem, sections 1--5} applies to any
complete finite-volume exterior. Its neighborhood is a neighborhood of
the \textbf{framed representation}, not only of the peripheral traces.
We recall the points preventing an exponentiated filling equation from
being mistaken for a metric.

On a compact truncated core choose a finite lifted partition of unity
\(\eta_{j,\gamma}\). In the future unit hyperboloid form \begin{equation}\label{mcg:eq-26}
Y_\rho(x)=\sum_{j,\gamma}\eta_{j,\gamma}(x)
                \rho(\gamma)\rho_0(\gamma)^{-1}x,\qquad
 E_\rho(x)=\frac{Y_\rho(x)}
                    {\sqrt{-\langle Y_\rho(x),Y_\rho(x)\rangle}} .
\end{equation} The future cone is convex; this is defined and equivariant. On a
compact fundamental part it is \(C^1\)-close to the identity, so is a
positive local diffeomorphism.

In a selected cusp frame the peripheral transformations are affine:
\(A_\rho(\lambda)z=a_\lambda z+b_\lambda\), on the original translation
lattice \(\Lambda\). A periodic partition of unity gives \begin{equation}\label{mcg:eq-27}
\begin{aligned}
d_\rho(w)&=\sum_{\lambda\in\Lambda}
   \eta(w-\lambda)A_\rho(\lambda)(w-\lambda),\\
F_\rho(w,t)&=(d_\rho(w),e^{t+\chi_\rho(w)}),
\end{aligned}
\end{equation} where \(\chi_\rho\) is the real-linear extension of
\(\lambda\mapsto\log|a_\lambda|\). These maps are exactly equivariant.
The first is a positive local diffeomorphism on a compact fundamental
parallelogram, and by equivariance everywhere. The second has positive
triangular differential on the entire tail. It can be joined to
\eqref{mcg:eq-26} by a \(C^1\)-small equivariant interpolation on a
fixed compact collar.

At the selected unit-charge critical point, the limiting version of
\eqref{mcg:eq-33} gives \(\operatorname{Im}\delta_e>0\) for either sign
of \(n_e\), since the limiting \(B_{e,N}\) in \eqref{mcg:eq-33} has positive real part. Consequently
\(\operatorname{Re}V_e=4\pi\operatorname{Im}\delta_e>0\), so the
longitude multiplier has \(|e^{V_e}|>1\). This is a nonunit
\textbf{modulus}, not merely the weaker assertion \(e^{V_e}\ne1\). The
commuting affine maps therefore have a second common fixed point
\(q=b/(1-a)\), using \(a=e^{V_e}\) for the longitude. As
\(\rho\to\rho_0\), \(q\) escapes to infinity, whereas the developed
fundamental parallelogram and its fixed generator paths remain bounded.
They miss \(q\); equivariance makes the entire development miss it. On
its simply connected cover set \(f=\log(d_\rho-q)\). The nonsingular
invariant metric \(|df|^2\) descends to a compact flat torus and is
complete. Consequently its development is the plane and its log periods
are an injective rank-two lattice.

For each fixed generator path of image length at most \(L\), in
\(|d_\rho|\le R\), its logarithmic increment is at most \(L/(|q|-R)\),
hence tends to zero. Its exponential is the corresponding similarity
multiplier. The difference from the small prescribed log \(U\) or \(V\)
is a multiple of \(2\pi i\) tending to zero, so is exactly zero. This
proves the actual log periods \textbf{before} multiplying them by
\(p_e\). One neighborhood works for all slopes.

After moving \(q\) to zero, the tail contains a full epigraph
\(y>A|d_\rho|\). With \[
 R_{\mathrm{tube}}=\operatorname{arsinh}(|d_\rho|/y),\quad
 \vartheta=\operatorname{Im}f,\quad
 X=\operatorname{Re}f+\tfrac12\log(1+y^2/|d_\rho|^2),
\] its metric becomes \begin{equation}\label{mcg:eq-28}
dR_{\mathrm{tube}}^2+\cosh^2(R_{\mathrm{tube}})\,dX^2
                   +\sinh^2(R_{\mathrm{tube}})\,d\vartheta^2 .
\end{equation} The primitive slope in \eqref{mcg:eq-22} has angle \(-2\pi\) and zero
axial period. A complementary primitive period has nonzero real part,
because the lattice has rank two. Adding \(R_{\mathrm{tube}}=0\)
therefore adds a smooth solid torus with precisely that meridian. Its
holonomy fixes both \(q,\infty\) and has multiplier one, so is the
identity; the representation descends to the filled group. This excludes
the parabolic-slope defect which a bare equation \(\exp(U+pV)=1\) would
permit.

At every main cusp all multipliers are one, so \eqref{mcg:eq-27} gives a
complete Euclidean torus link and a full horoball quotient. The
resulting metric is complete and finite-volume. Its universal developing
local isometry covers \(\mathbb H^3\), proving discrete faithful
holonomy.

The representation volume agrees with this filled metric's volume using
the same finite boundary triangulation: compare its cusp map with the
axis-valued map by the common conical boundary homotopy, whose volume
tends to zero, and extend that triangulation through the solid torus
with all vertices on the core axis, so every added tetrahedron has zero
volume. This is the filling-chain comparison in
\href{https://sashakolpakov.github.io/murakamis_dream/notes/general_source_regulator_hyperbolic_equality.html}{the
general regulator theorem, section 6}; it retains the same oriented
fundamental class.

Our family satisfies the neighborhood hypothesis: \eqref{mcg:eq-8},
\eqref{mcg:eq-14}, \eqref{mcg:eq-16}, and their chosen flags are
holomorphic at zero and converge to their original complete matrices as
all \(|n_e|\to\infty\). The first-charge solution above therefore lies
in this fixed neighborhood after one diagram-dependent threshold.
Inserting \eqref{mcg:eq-25} now gives the volume of the actual filled
manifold.

Finally, filling a crossing circle along its recorded
\(\mu_e+p_e\lambda_e\) is the local surgery in the original
twice-punctured disk collar which inserts \(-2p_e\) crossings. The disks
and collars are those constructed in \(S^3\) in
section~\ref{mcg:sec-2}. Thus the filled topology is the
original diagram with powers \eqref{mcg:eq-1}, not a comparison diagram,
a tree replacement, or an unnamed representation component.

\subsection{Analytic interface and the remaining full-source
problem}\label{mcg:sec-8}

The independent analytic verification is recorded in
\hyperref[mca:unit]{the
analytic packet and whole-period proof}. Its inputs are the exact source
\eqref{mcg:eq-5}--\eqref{mcg:eq-7}, the positive complete radii and the
original primitive face lattice. The relevant graph facts hold here:
deleting the protected adjacent vertices from a simple triangulated
sphere leaves its surviving ordinary graph connected; a matching outer
pair is not the protected pair; every matching factor has \(x+y>0\) and
at least one nonmonomial triangle. These give the ordinary Hessian rank
and primitive angular differences used by the complex coefficient
saddle. They do not require a flat gap forest.

The companion proof establishes the local holomorphic packet and global
complex-strip spectral gap, including all odd simple-node normals. It
then applies the original contour proof to \eqref{mcg:eq-7}, with real
spectators during infinite-tail removal, exact endpoint currents, and
independent signed reciprocal slopes. That proof returns the nonzero
first-charge asymptotic, not just an upper bound on its action.
Section~\ref{mcg:sec-9} records its exact normalization.

For an unchanged opposite-offset even bipyramid, its displayed circle
packing is a positive certificate and its oriented top and bottom
quartets are exactly the general Morse source \eqref{mcg:eq-5}.
Consequently this theorem applies to its all-central term at every fixed
size, with the original parity/knot qualification. No Conway-tree
decomposition or new plat presentation is required.

What remains for its \textbf{full} invariant is a strict estimate of the
signed contribution outside central, for example \begin{equation}\label{mcg:eq-29}
\bigl|\mathcal I_{N,\mathrm{full}}
            -\mathcal I_{N,\mathrm C}\bigr|
 \le \exp\!\left[
       N\left\{\frac{\operatorname{Vol}(S^3\setminus K_{\boldsymbol p})}
                       {2\pi}-\eta\right\}+o(N)\right].
\end{equation} Failure of complete-radius balance at every mark neither invalidates
the all-central theorem nor proves \eqref{mcg:eq-29} false. It removes one
known sufficient method of proving \eqref{mcg:eq-29}. The present result
makes no claim of cancellation or domination for those remaining
sectors.

\subsection{The nonzero asymptotic with its finite
normalization}\label{mcg:sec-9}

Here is the analytic statement actually joined to the geometry, with the
finite quantities retained. On one fixed complex polydisk there are even
holomorphic face-critical actions \(A_N\) and face Hessians \(H_N\),
converging with every fixed derivative to their limiting quantities,
such that \begin{equation}\label{mcg:eq-30}
C_N^{\mathrm{raw}}(\boldsymbol\delta)
 =\frac{(-1)^{dR}N^{2R}\exp(NA_N(\boldsymbol\delta))}
 {(2\pi N)^{(R-1)/2}\sqrt{\det H_N(\boldsymbol\delta)}}
                      (1+O(N^{-1})).
\end{equation} The action includes all finite anchors, sine and beta terms and any
\((\log N)/N\) terms; only the displayed \(N^{2R}\) was removed. Its
limit is \eqref{mcg:eq-17}. The square root in \eqref{mcg:eq-30} is
continued from the positive complete face Hessian.

For every fixed \(\eta>0\), there are \(\gamma_\eta,\tau_\eta>0\) such
that, on the entire real spectral cube and the indicated thin strip, \begin{equation}\label{mcg:eq-31}
\begin{gathered}
 |C_N^{\mathrm{raw}}(\boldsymbol\delta)|
       \le e^{N(\Phi_D-\gamma_\eta)},\\
 \max_e|\operatorname{Re}\delta_e|\ge\eta,\quad
 \operatorname{Re}\delta_e\in[-1/2,1/2],\quad
 |\operatorname{Im}\delta_e|\le\tau_\eta .
\end{gathered}
\end{equation} Fixed derivatives satisfy the same type of bound on smaller domains.
The local cofactor interpolation retains the odd simple-node derivative,
while the shared angular estimate uses the connected ordinary graph and
its primitive face-difference rows. The separated real-spectral estimate
comes from the exact cofactor starting-row maximum. Its only monomial
exception is harmless because every crossing factor has a nonmonomial
triangle. These prove \eqref{mcg:eq-30},\eqref{mcg:eq-31} for the all-central projection of the exact
coefficient \eqref{mcg:eq-5}; they are not inferred from a value-only
local recurrence.

Let \(\boldsymbol\delta_{N,\boldsymbol1}\) be the finite critical point
of \[
 A_N(\boldsymbol\delta)
       +\pi i\sum_e n_e\delta_e^2+2\pi i\sum_e\delta_e ,
\] let \(S_{N,\boldsymbol1}\) be its critical value, and let
\(Q_{N,\boldsymbol1}\) be its negative spectral Hessian. Then the
\textbf{entire} period \eqref{mcg:eq-7} is \begin{equation}\label{mcg:eq-32}
\begin{split}
\mathcal I_{N,\mathrm C}
={}&q^{d^2\sum_e n_e/4}(-2i)^R\sqrt{2\pi}\,N^{R+1/2}\\
&\ \cdot
\frac{\prod_e\sin(2\pi\delta_{N,\boldsymbol1,e})}
 {\sqrt{\det H_N(\boldsymbol\delta_{N,\boldsymbol1})}
      \sqrt{\det Q_{N,\boldsymbol1}}}
\exp(NS_{N,\boldsymbol1})\\
&\ \cdot
\left(1+O_{\boldsymbol n}(N^{-1})
                +O_{\boldsymbol n}(e^{-\kappa_{\boldsymbol n}N})\right).
\end{split}
\end{equation} The spectral determinant uses the Gaussian branch on the cone with
positive real part. The sine and both determinant factors have a nonzero
limiting product for fixed sufficiently long powers.

For clarity, the charge proof has four logically separate parts.
Separate evenness writes
\(\partial_eA_N=-2\delta_e B_{e,N}(\boldsymbol\delta^2)\), where
\(B_{e,N}\to C_e>0\). The critical equations for positive charge
magnitudes give \begin{equation}\label{mcg:eq-33}
\delta_e=-\frac{\ell_e/n_e}
       {1+iB_{e,N}(\boldsymbol\delta^2)/(\pi n_e)},\qquad
 \operatorname{Im}\delta_e
  =\frac{\ell_e}{\pi n_e^2}
      \frac{\operatorname{Re}B_{e,N}}
          {|1+iB_{e,N}/(\pi n_e)|^2}>0.
\end{equation} An actual outer local cube is shifted to this saddle, with every side
current retained below the saddle value. Its negative Hessian has
positive real part and its Gaussian coefficient is one. The face and
spectral Gaussian powers combine to \begin{equation}\label{mcg:eq-34}
N^{2R}(2\pi N)^{-(R-1)/2}(2\pi/N)^{R/2}
                         =\sqrt{2\pi}\,N^{R+1/2}.
\end{equation}

Next, differentiating the critical value gives the joint comparison \begin{equation}\label{mcg:eq-35}
\operatorname{Re}S_{N,\boldsymbol1}
       -\operatorname{Re}S_{N,\boldsymbol\ell}
   \ge\sum_e\frac{C_e}{2n_e^2}(\ell_e^2-1)
\end{equation} on a fixed scaled local charge box. This is why unequal slopes cause
no loss of the theorem: it compares the critical values in one joint
family, rather than discarding local errors independently.

Third, infinite far charges are removed on whole real intervals, with
the other coordinates kept as real spectators. The sine factor gives
zero endpoint values and equal first derivatives; three integrations by
parts retain the second-derivative endpoint current. On the shifted
contour the two vertical endpoints are paired before estimation. Their
hyperbolic-sine difference and \eqref{mcg:eq-31} give a summable strict
gap. The degree-\(d\) cosine bound controls the remaining horizontal
interval. Only then is the remaining finite charge set divided into
local and bounded-bad charges. A separate smaller cube, not the outer
saddle cube, controls the latter by a one-coordinate shift. All
connectors and all original endpoints are included.

Finally, reflecting any charge sign changes both its sine integral and
its charge weight, so all \(2^R\) unit signs agree. The raw parity sign
cancels exactly: \begin{equation}\label{mcg:eq-36}
(-1)^{dR}(-1)^{(N+1)R}=1 .
\end{equation} This leaves the original \(q\) phase and the factor \((-2i)^R\) in
\eqref{mcg:eq-32}. Neither the sign choice of a power nor the parity of
\(N\) can force cancellation of that leading coefficient.

Equations \eqref{mcg:eq-30}--\eqref{mcg:eq-36} give the nonzero lower as
well as the upper asymptotic. Taking its logarithm and using
\eqref{mcg:eq-25} and the metric construction proves
\eqref{mcg:eq-2},\eqref{mcg:eq-3}. Thus the flat-plat wording in the
older all-central theorem is a limitation of its stated input, not a
remaining mathematical restriction after these separate Morse source,
analytic, and geometric arguments.
